% Options for packages loaded elsewhere
\PassOptionsToPackage{unicode}{hyperref}
\PassOptionsToPackage{hyphens}{url}
\PassOptionsToPackage{dvipsnames,svgnames,x11names}{xcolor}
\documentclass[
  11pt,
]{article}
\usepackage{xcolor}
\usepackage[margin=27mm]{geometry}
\usepackage{amsmath,amssymb}
\setcounter{secnumdepth}{-\maxdimen} % remove section numbering
\usepackage{iftex}
\ifPDFTeX
  \usepackage[T1]{fontenc}
  \usepackage[utf8]{inputenc}
  \usepackage{textcomp} % provide euro and other symbols
\else % if luatex or xetex
  \usepackage{unicode-math} % this also loads fontspec
  \defaultfontfeatures{Scale=MatchLowercase}
  \defaultfontfeatures[\rmfamily]{Ligatures=TeX,Scale=1}
\fi
\usepackage{lmodern}
\ifPDFTeX\else
  % xetex/luatex font selection
  \setmainfont[]{Cambria}
  \setmonofont[]{Consolas}
\fi
% Use upquote if available, for straight quotes in verbatim environments
\IfFileExists{upquote.sty}{\usepackage{upquote}}{}
\IfFileExists{microtype.sty}{% use microtype if available
  \usepackage[]{microtype}
  \UseMicrotypeSet[protrusion]{basicmath} % disable protrusion for tt fonts
}{}
\makeatletter
\@ifundefined{KOMAClassName}{% if non-KOMA class
  \IfFileExists{parskip.sty}{%
    \usepackage{parskip}
  }{% else
    \setlength{\parindent}{0pt}
    \setlength{\parskip}{6pt plus 2pt minus 1pt}}
}{% if KOMA class
  \KOMAoptions{parskip=half}}
\makeatother
\setlength{\emergencystretch}{3em} % prevent overfull lines
\providecommand{\tightlist}{%
  \setlength{\itemsep}{0pt}\setlength{\parskip}{0pt}}
\usepackage{bookmark}
\IfFileExists{xurl.sty}{\usepackage{xurl}}{} % add URL line breaks if available
\urlstyle{same}
\hypersetup{
  pdftitle={WEIS-specific prior-art map for the seventh-capital thesis},
  pdfauthor={The Privacy-is-Value Research Programme},
  colorlinks=true,
  linkcolor={black},
  filecolor={Maroon},
  citecolor={Blue},
  urlcolor={blue},
  pdfcreator={LaTeX via pandoc}}

\title{WEIS-specific prior-art map for the seventh-capital thesis}
\usepackage{etoolbox}
\makeatletter
\providecommand{\subtitle}[1]{% add subtitle to \maketitle
  \apptocmd{\@title}{\par {\large #1 \par}}{}{}
}
\makeatother
\subtitle{WP-14 · tier A · WEIS-2027 · 2026-07-14}
\author{The Privacy-is-Value Research Programme}
\date{2026-07-14}

\begin{document}
\maketitle

\section{WEIS Prior-Art Map --- WP-14 (seventh capital / S--M separation
/
priced-elsewhere)}\label{weis-prior-art-map-wp-14-seventh-capital-sm-separation-priced-elsewhere}

\textbf{Role discipline (A10):} external, resolvable citations only.
Every entry below was web-verified for author, title, year and venue
against a primary or near-primary record (archive PDF, publisher landing
page, or author page). Entries I could not fully confirm are marked
\textbf{UNVERIFIED}. No citation is invented. Vocabulary follows GR-4:
the observing agent is the \textbf{boundary agent S}, the acting agent
is the \textbf{delegation agent M}.

\textbf{Our thesis, stated for positioning (WP-14):} 1.
behavioural/personal data is a distinct capital class (the ``seventh
capital''); 2. architectural separation between an observing
\textbf{boundary agent S} and an acting \textbf{delegation agent M}
yields information-theoretic (not policy) guarantees, time-indexed as
R(t); 3. privacy has a market price that is currently captured by a
party other than the data subject.

Relation tags used: \textbf{AGREES} · \textbf{EXTENDS-WE-BUILD-ON} ·
\textbf{CONTRADICTS-WE-MUST-ANSWER} · \textbf{GAP-WE-FILL}.

\begin{center}\rule{0.5\linewidth}{0.5pt}\end{center}

\subsection{Coverage statement (no silent
gaps)}\label{coverage-statement-no-silent-gaps}

WEIS per-year sites use two live patterns:
\texttt{weis\textless{}YEAR\textgreater{}.econinfosec.org/}
(accepted-papers/program pages) and
\texttt{econinfosec.org/archive/weis\textless{}YEAR\textgreater{}/papers/\textless{}n\textgreater{}.pdf}
(per-paper PDFs). Findings:

\begin{itemize}
\tightlist
\item
  \textbf{Fully enumerated accepted-paper lists (paper-by-paper):} WEIS
  2022 and WEIS 2023. WEIS 2023's full accepted list contains
  \textbf{no} privacy-valuation / personal-data-market paper
  (nearest-adjacent: Cecere/Jean/Lefrere/Tucker on platform regulatory
  compliance). WEIS 2022 contains the Collis et al.~valuation paper
  (entry P3) plus Telang/Garg on app privacy-label demand (adjacent, not
  valuation).
\item
  \textbf{Confirmed archive PDF hosting (individual papers retrieved or
  path-verified):} WEIS 2007 (entry P1), WEIS 2015 (entry P2).
\item
  \textbf{NOT enumerated paper-by-paper:} WEIS 2002--2006, 2008--2014,
  2016--2021, 2024, 2025. For these years I searched by author and topic
  rather than walking every accepted list, so a non-flagship
  privacy-economics paper in those years could exist that this map does
  not name. The flagship privacy-valuation lineage (Acquisti and
  collaborators; the privacy-paradox field experiments) is captured
  across the whole span via targeted search.
\item
  \textbf{Tooling limitation (not a coverage gap):} the WebFetch
  renderer cannot parse PDF binaries in this environment (no poppler),
  so numeric results were taken from readable HTML mirrors, publisher
  landing pages, author-hosted abstracts, and press summaries rather
  than from the PDF body directly. Where a number rests on a secondary
  summary rather than the paper's own table, the entry says so.
\item
  \textbf{Archive index page} \texttt{econinfosec.org/weis-archive/}
  lists editions 2002--2023 but carries no per-paper metadata; it
  returned 403 on the \texttt{/archive/weis2015/} directory index even
  though individual \texttt{/archive/weis2015/papers/*.pdf} files
  resolve.
\end{itemize}

\begin{center}\rule{0.5\linewidth}{0.5pt}\end{center}

\subsection{Verified entries}\label{verified-entries}

\subsubsection{P1 --- Grossklags \& Acquisti, WEIS
2007}\label{p1-grossklags-acquisti-weis-2007}

\begin{itemize}
\tightlist
\item
  \textbf{CITATION:} Jens Grossklags, Alessandro Acquisti. ``When 25
  Cents Is Too Much: An Experiment on Willingness-To-Sell and
  Willingness-To-Protect Personal Information.'' \emph{6th Workshop on
  the Economics of Information Security (WEIS 2007)}, Pittsburgh, June
  2007. Archive PDF:
  \texttt{https://econinfosec.org/archive/weis2007/papers/66.pdf} (also
  mirrored at
  \texttt{cs.cit.tum.de/.../2007-WEIS-Grossklags-Acquisti.pdf}).
  VERIFIED (author, title, year, WEIS venue all confirmed across archive
  path, TUM mirror, Semantic Scholar, DBLP/researchr).
\item
  \textbf{FINDING:} Yes/no offers and open valuations show a strong
  majority preferring money over data even when the monetary advantage
  of releasing (or not protecting) is very small; measured
  willingness-to-protect (a WTP) is far below
  willingness-to-accept-to-sell (a WTA).
\item
  \textbf{RELATION:} EXTENDS-WE-BUILD-ON --- this is the WEIS-native
  root of ``subjects under-price their own data,'' which is the
  demand-side half of our ``price captured by someone else.'' It also
  sets up a \textbf{CONTRADICTS-WE-MUST-ANSWER} tension: if subjects
  value data at cents, a referee will ask how it can be a distinct
  capital class. Our answer must separate individual WTP from the
  aggregate market price.
\item
  \textbf{HANDLE FOR VALUATION:} the paper's own threshold framing is
  the ``25 cents'' of the title (a per-decision offer at which subjects
  still sold). Exact table figures were not extracted from the PDF body
  (binary-render limit); treat the 25-cent figure as the load-bearing
  quotable and pull precise means/medians from the PDF in the valuation
  rebuild.
\end{itemize}

\subsubsection{P2 --- Cofone, WEIS 2015}\label{p2-cofone-weis-2015}

\begin{itemize}
\tightlist
\item
  \textbf{CITATION:} Ignacio N. Cofone. ``The Value of Privacy: Keeping
  the Money Where the Mouth Is.'' \emph{WEIS 2015}, Delft. Archive PDF:
  \texttt{https://econinfosec.org/archive/weis2015/papers/WEIS\_2015\_cofone.pdf}.
  VERIFIED (title, author, WEIS 2015 venue via archive path and Semantic
  Scholar).
\item
  \textbf{FINDING:} Argues the privacy paradox is explicable inside a
  rational-choice framework; low willingness-to-pay for privacy explains
  the thin market for privacy-enhancing technologies (uses Tor's
  free-but-used status and Facebook as case studies).
\item
  \textbf{RELATION:} CONTRADICTS-WE-MUST-ANSWER --- if low WTP is
  rational rather than a market failure, an architectural remedy looks
  unnecessary. We must answer by distinguishing a rational low WTP for
  \emph{policy} privacy from an information-theoretic leakage bound R(t)
  that a rational subject cannot purchase at any price under the current
  architecture. Partial \textbf{AGREES} on ``someone else captures the
  value.''
\item
  \textbf{HANDLE FOR VALUATION:} cites Facebook as \textasciitilde850
  million users generating a firm valuation of \textasciitilde US\$100
  billion (2015 framing) as an illustration of subject-generated value
  accruing to the platform. This is a scene-setting figure from a
  secondary case description, not an estimated per-user price; use only
  as context.
\end{itemize}

\subsubsection{P3 --- Collis, Moehring, Sen \& Acquisti, WEIS 2022 ★
strongest empirical
anchor}\label{p3-collis-moehring-sen-acquisti-weis-2022-strongest-empirical-anchor}

\begin{itemize}
\tightlist
\item
  \textbf{CITATION:} Avinash Collis, Alex Moehring, Ananya Sen,
  Alessandro Acquisti. ``Information Frictions and Heterogeneity in
  Valuations of Personal Data.'' \emph{WEIS 2022}. PDF:
  \texttt{https://weis2022.econinfosec.org/wp-content/uploads/sites/10/2022/06/weis22-collis.pdf};
  SSRN 3974826 (``\ldots Valuations of Social Media Data'', Nov 2021).
  VERIFIED (accepted-papers list, PDF URL, SSRN, CMU/Penn summaries).
\item
  \textbf{FINDING:} Incentive-compatible mechanism elicits how much
  compensation users require to share their Facebook data; finds large
  dispersion and systematic under-valuation by women, Black, and
  low-income users; an information treatment compresses the dispersion.
\item
  \textbf{RELATION:} EXTENDS-WE-BUILD-ON (empirics) and
  \textbf{GAP-WE-FILL} (mechanism). Their remedy for the frictions is
  \emph{information provision} --- a policy lever. Ours is
  \emph{architecture} (S--M separation). Their demographic dispersion is
  direct evidence that consent/notice regimes misprice, which our
  information-theoretic guarantee sidesteps.
\item
  \textbf{HANDLE FOR VALUATION (quote precisely):} median WTA to share
  \textbf{all} Facebook data = \textbf{US\$750} (YouGov sample) and
  \textbf{US\$1,000} (DDP sample), one-time (not monthly), 2021
  elicitation. Demographic medians (YouGov): Black \textbf{\$500} vs
  White \textbf{\$1,000}; Female \textbf{\$600} vs Male
  \textbf{\$1,000}. (Figures via the Penn Carey Law summary of the
  paper; confirm against the PDF's tables in the rebuild.)
\end{itemize}

\subsubsection{P4 --- Acquisti, John \& Loewenstein, ``What Is Privacy
Worth?''
(WEIS-adjacent)}\label{p4-acquisti-john-loewenstein-what-is-privacy-worth-weis-adjacent}

\begin{itemize}
\tightlist
\item
  \textbf{CITATION:} Alessandro Acquisti, Leslie K. John, George
  Loewenstein. ``What Is Privacy Worth?'' \emph{Journal of Legal
  Studies} 42(2), Article 1, 2013 (DOI 10.1086/671754). Earlier
  presented as a workshop paper (WISE 2009 version at
  \texttt{pages.stern.nyu.edu/\textasciitilde{}bakos/wise/papers/wise2009-6a1\_paper.pdf};
  author copy at heinz.cmu.edu). VERIFIED (JLS volume/issue, three
  authors, year via ideas.repec, chicagounbound, author page). Venue is
  JLS, not WEIS --- labelled adjacent.
\item
  \textbf{FINDING:} Field experiment with a \$10 anonymous vs \$12
  identified gift-card design shows large endowment and order effects:
  subjects endowed with anonymity demand far more to give it up (WTA)
  than subjects without it will pay to obtain it (WTP); valuations
  cluster at focal/extreme points.
\item
  \textbf{RELATION:} EXTENDS-WE-BUILD-ON --- instability and
  framing-dependence of privacy valuations undercut consent/notice
  (policy) guarantees and motivate a guarantee that does not depend on
  subject valuation at all, i.e.~our information-theoretic R(t).
\item
  \textbf{HANDLE FOR VALUATION:} design constants are exact and
  quotable: \textbf{\$10 anonymous card} vs \textbf{\$12 identified
  card} (a \textbf{\$2} wedge to switch); WTA:WTP disparity reported as
  roughly \textbf{5×} (secondary summaries; verify the exact ratio in
  the JLS tables for the rebuild).
\end{itemize}

\subsubsection{P5 --- Beresford, Kübler \& Preibusch, ``Unwillingness to
Pay for Privacy''
(adjacent)}\label{p5-beresford-kuxfcbler-preibusch-unwillingness-to-pay-for-privacy-adjacent}

\begin{itemize}
\tightlist
\item
  \textbf{CITATION:} Alastair R. Beresford, Dorothea Kübler, Sören
  Preibusch. ``Unwillingness to Pay for Privacy: A Field Experiment.''
  \emph{Economics Letters} 117(1): 25--27, 2012; earlier IZA Discussion
  Paper 5017 (2010) and WZB working paper. VERIFIED (IZA, SSRN 1634484,
  ScienceDirect, WZB PDF). Venue is Economics Letters / IZA, not WEIS
  --- adjacent, but the canonical privacy-paradox field experiment.
\item
  \textbf{FINDING:} Subjects buy a DVD from one of two
  otherwise-identical stores differing only in how much personal data
  they demand. With a \textbf{€1} discount at the data-hungry store,
  almost everyone gives up the extra data; with equal prices, purchases
  split roughly evenly despite stated privacy concern.
\item
  \textbf{RELATION:} CONTRADICTS-WE-MUST-ANSWER --- the hardest
  empirical for our capital-class claim (``people sell everything for
  €1''). We answer: market price ≠ subject WTP; the seventh-capital
  claim concerns the aggregate value captured downstream, not the
  subject's marginal willingness to pay at the point of sale.
\item
  \textbf{HANDLE FOR VALUATION:} the \textbf{€1} price wedge that flips
  nearly all buyers to the data-demanding store (2010--2012). Quotable
  as an upper bound on point-of-sale subject WTP.
\end{itemize}

\subsubsection{P6 --- Acquisti, Taylor \& Wagman, ``The Economics of
Privacy'' (adjacent survey ---
must-cite)}\label{p6-acquisti-taylor-wagman-the-economics-of-privacy-adjacent-survey-must-cite}

\begin{itemize}
\tightlist
\item
  \textbf{CITATION:} Alessandro Acquisti, Curtis R. Taylor, Liad Wagman.
  ``The Economics of Privacy.'' \emph{Journal of Economic Literature}
  54(2): 442--492, 2016 (DOI 10.1257/jel.54.2.442). VERIFIED (AEA,
  ideas.repec, author PDF, Duke Scholars). Adjacent venue (JEL) but
  authored by WEIS-community principals and is THE field survey.
\item
  \textbf{FINDING:} Canonical survey of the economic value and
  consequences of protecting vs disclosing personal information, the
  trade-offs, externalities, and consumer decision-making anomalies.
\item
  \textbf{RELATION:} GAP-WE-FILL + EXTENDS-WE-BUILD-ON --- the survey
  treats personal data as an economic good with externalities and
  information asymmetries; it does \textbf{not} frame data as a distinct
  capital class, and it contains no architectural or
  information-theoretic (R(t)) treatment. This is the map against which
  our novelty is measured; a referee will expect us to locate ourselves
  inside it.
\item
  \textbf{HANDLE FOR VALUATION:} survey, not a primary estimate. Use it
  to source the accepted range of WTP/WTA estimates rather than as a
  single number.
\end{itemize}

\subsubsection{P7 --- Acquisti \& Grossklags, ``Privacy and Rationality
in Individual Decision Making'' (WEIS 2004
origin)}\label{p7-acquisti-grossklags-privacy-and-rationality-in-individual-decision-making-weis-2004-origin}

\begin{itemize}
\tightlist
\item
  \textbf{CITATION:} Alessandro Acquisti, Jens Grossklags. ``Privacy and
  Rationality in Individual Decision Making.'' \emph{IEEE Security \&
  Privacy} 3(1): 26--33, 2005 (DOI 10.1109/MSP.2005.22). An earlier
  version was presented at WEIS 2004. VERIFIED for the IEEE S\&P
  publication (IEEE/ACM DL, PSU, Semantic Scholar); the WEIS 2004
  presentation attribution is asserted by SSRN/author records --- mark
  the WEIS-2004 line \textbf{UNVERIFIED} and cite the IEEE venue as
  primary.
\item
  \textbf{FINDING:} Empirically and theoretically documents bounded
  rationality and information asymmetry in privacy decisions; subjects
  behave inconsistently with rational-agent models.
\item
  \textbf{RELATION:} EXTENDS-WE-BUILD-ON --- foundational reason to
  replace consent/notice (which presumes a rational, informed subject)
  with an architectural guarantee that holds regardless of subject
  rationality.
\item
  \textbf{HANDLE FOR VALUATION:} methodology constants only: a
  \textbf{US\$16} lump-sum for an online survey, \textbf{119}
  respondents (May 2004). Not a data-value estimate.
\end{itemize}

\subsubsection{P8 --- Arrieta-Ibarra, Goff, Jiménez-Hernández, Lanier \&
Weyl, ``Should We Treat Data as Labor?'' (adjacent --- rival framing) ★
key positioning
fight}\label{p8-arrieta-ibarra-goff-jimuxe9nez-hernuxe1ndez-lanier-weyl-should-we-treat-data-as-labor-adjacent-rival-framing-key-positioning-fight}

\begin{itemize}
\tightlist
\item
  \textbf{CITATION:} Imanol Arrieta-Ibarra, Leonard Goff, Diego
  Jiménez-Hernández, Jaron Lanier, E. Glen Weyl. ``Should We Treat Data
  as Labor? Moving Beyond `Free'.'' \emph{AEA Papers and Proceedings}
  108: 38--42, 2018 (DOI 10.1257/pandp.20181003). VERIFIED (AEA, SSRN
  3093683, EconPapers vol/pages).
\item
  \textbf{FINDING:} Argues data currently treated as free capital should
  instead be treated (at least partly) as \textbf{labour} --- a
  compensated factor input --- to restore a functioning market for user
  contributions.
\item
  \textbf{RELATION:} CONTRADICTS-WE-MUST-ANSWER --- the single closest
  rival to our framing. They say data is \emph{labour}; we say it is a
  distinct \emph{capital} class (``seventh capital''). Our related-work
  section must draw the labour-vs-capital line explicitly and say why an
  information-theoretic boundary (S--M) changes the accounting in a way
  a labour-compensation scheme does not.
\item
  \textbf{HANDLE FOR VALUATION:} conceptual, no primary estimate.
\end{itemize}

\subsubsection{P9 --- Posner \& Weyl, ``Data as Labor'' (Radical
Markets) (adjacent --- book form of
P8)}\label{p9-posner-weyl-data-as-labor-radical-markets-adjacent-book-form-of-p8}

\begin{itemize}
\tightlist
\item
  \textbf{CITATION:} Eric A. Posner, E. Glen Weyl. ``Data as Labor,''
  chapter 5 in \emph{Radical Markets: Uprooting Capitalism and Democracy
  for a Just Society}. Princeton University Press, 2018. VERIFIED (De
  Gruyter/Princeton chapter listing, radicalmarkets.com). Book, not a
  venue paper.
\item
  \textbf{FINDING:} Full-length statement of the data-as-labour
  programme, including data-union / data-cooperative institutions as the
  labour-market analogue.
\item
  \textbf{RELATION:} CONTRADICTS-WE-MUST-ANSWER --- same axis as P8;
  cite the book once as the canonical long-form and engage P8 as the
  citable venue paper.
\item
  \textbf{HANDLE FOR VALUATION:} conceptual.
\end{itemize}

\begin{center}\rule{0.5\linewidth}{0.5pt}\end{center}

\subsection{Adjacent / watch-list (named, not yet
load-bearing)}\label{adjacent-watch-list-named-not-yet-load-bearing}

\begin{itemize}
\tightlist
\item
  \textbf{Cecere, Jean, Lefrere \& Tucker}, ``Trade-offs in Automating
  Platform Regulatory Compliance By Algorithm,'' WEIS 2023 --- the only
  2023 paper touching platform data governance; not valuation. VERIFIED
  (2023 accepted list).
\item
  \textbf{Telang \& Garg}, ``Impact of App Privacy Label Disclosure on
  Demand: An Empirical Analysis,'' WEIS 2022 --- privacy-label demand
  effects, adjacent to consent/notice economics; not a valuation
  estimate. VERIFIED (2022 accepted list).
\item
  \textbf{Kuerbis \& Mueller}, ``Exploring the role of data enclosure in
  the digital political economy,'' WEIS 2022 --- ``data enclosure'' is
  conceptually near the capital-class claim; worth reading before final
  related-work. VERIFIED (2022 accepted list).
\item
  \textbf{Spiekermann \& Korunovska}, ``Towards a value theory for
  personal data,'' \emph{Journal of Information Technology} 2017 ---
  survey-experiment value theory (1,269 Facebook users); adjacent,
  potential valuation input. UNVERIFIED as to exact figures (only
  landing page seen).
\end{itemize}

\begin{center}\rule{0.5\linewidth}{0.5pt}\end{center}

\subsection{SYNTHESIS}\label{synthesis}

\subsubsection{(1) Where the seventh-capital thesis is genuinely novel
vs.~what WEIS has already
published}\label{where-the-seventh-capital-thesis-is-genuinely-novel-vs.-what-weis-has-already-published}

WEIS and its adjacent literature have thoroughly established three
things our paper leans on but must \textbf{not} re-claim as novel: (a)
subjects systematically under-price their own data and behave
irrationally about it (P1, P7); (b) privacy valuations are unstable,
endowment- and framing-dependent (P4); and (c) the WTA/WTP gap and the
privacy paradox are real and measurable (P1, P4, P5). Two things are
contested rather than settled and are where we must fight: whether the
paradox implies a market failure at all (Cofone, P2, says no), and
whether data is best modelled as \textbf{labour} (Arrieta-Ibarra/Weyl,
P8/P9) or as something else.

Genuinely novel, and absent from the WEIS body of work as surveyed: -
\textbf{Data as a distinct capital class with its own accounting}
(``seventh capital''). The field's default is data-as-economic-good (P6)
or data-as-labour (P8/P9). No WEIS paper found frames behavioural data
as a seventh capital class. - \textbf{An information-theoretic, not
policy, guarantee.} Every WEIS/adjacent treatment found operates on
incentives, disclosure, consent, notice, or compensation --- all policy
instruments. None derives a leakage bound from the \emph{architecture}
of the observing/acting split (our boundary agent S vs delegation agent
M). - \textbf{Time-indexed R(t).} The literature reports static
valuations and static paradoxes; a time-indexed reconstruction bound
tied to an architectural separation has no precedent in the corpus
mapped here.

The ``priced-elsewhere'' claim (our point 3) is the \emph{least} novel:
P1, P2, P3, P5 all already show value flowing away from the subject. We
should present point 3 as well-supported established ground, and spend
our novelty budget on points 1 (capital class) and 2
(architectural/information-theoretic guarantee).

\subsubsection{(2) The WEIS (+ must-cite adjacent) papers our
related-work MUST engage or a referee will reject
us}\label{the-weis-must-cite-adjacent-papers-our-related-work-must-engage-or-a-referee-will-reject-us}

\begin{enumerate}
\def\labelenumi{\arabic{enumi}.}
\tightlist
\item
  \textbf{Acquisti, Taylor \& Wagman, JEL 2016 (P6)} --- the field
  survey; not citing it reads as not knowing the field.
\item
  \textbf{Collis, Moehring, Sen \& Acquisti, WEIS 2022 (P3)} --- the
  current WEIS-native empirical on data valuation; our valuation section
  must reconcile with it.
\item
  \textbf{Arrieta-Ibarra, Goff, Jiménez-Hernández, Lanier \& Weyl, 2018
  (P8)} --- the rival capital-vs-labour framing; we must draw the line
  explicitly.
\item
  \textbf{Acquisti, John \& Loewenstein, JLS 2013 (P4)} --- the
  canonical WTA/WTP-instability result; grounds our ``policy guarantees
  are weak'' move.
\item
  \textbf{Grossklags \& Acquisti, WEIS 2007 (P1)} --- the WEIS-native
  root of subject under-pricing; positions us inside WEIS's own lineage
  rather than beside it.
\end{enumerate}

(Beresford et al., P5, is the sixth and is the specific paper a
sceptical referee will wield against the capital-class claim; pre-empt
it.)

\subsubsection{(3) Strongest empirical numbers --- shortlist for the
valuation rebuild
(A12)}\label{strongest-empirical-numbers-shortlist-for-the-valuation-rebuild-a12}

Ranked by strength and directness for a per-subject data-price estimate:
1. \textbf{Collis et al.~2022 (P3):} median WTA \textbf{\$750 / \$1,000}
(one-time, all Facebook data, 2021), with demographic splits (\$500
Black / \$1,000 White; \$600 female / \$1,000 male). Strongest, most
recent, WEIS-native. Verify against the PDF tables. 2. \textbf{Acquisti,
John \& Loewenstein 2013 (P4):} \textbf{\$10 anonymous vs \$12
identified} gift-card design; \textbf{\$2} switching wedge;
\textasciitilde{}\textbf{5×} WTA:WTP disparity. Exact design constants;
ratio to be confirmed in JLS tables. 3. \textbf{Beresford et al.~2012
(P5):} \textbf{€1} discount flips nearly all buyers to surrender extra
personal data --- an upper bound on point-of-sale subject WTP. 4.
\textbf{Grossklags \& Acquisti 2007 (P1):} the \textbf{25-cent}
sell-threshold framing; extract precise means/medians from the PDF in
the rebuild. 5. \textbf{Cofone 2015 (P2):} \textasciitilde{}\textbf{850M
users / \textasciitilde\$100B} Facebook valuation --- context only, a
firm-value illustration, not a per-user estimate.

\textbf{Caveat for A12:} figures 1, 2 and 4 currently rest partly on
secondary summaries (press releases, publisher landing pages, law-school
blog) because the PDF bodies did not render in this environment. Each
must be re-confirmed against the paper's own tables before it enters a
TIER-A/TIER-S valuation figure. Do not let any of these numbers cross
the GR-3 fiat-value fence: they are external literature quotations for
positioning, not canon value claims.

\end{document}
